% ============================================================================
% Soluciones del Módulo 3 - Visualización de datos con ggplot2
% ============================================================================

\begin{solucion}{m3-1}
Contamos los tickets por tienda con \codigo{count()} y graficamos con
\codigo{geom\_col()} en horizontal, ordenando con \codigo{fct\_reorder()}.
\end{solucion}

\begin{rcode}
library(tidyverse)
library(scales)
color_principal <- "#2a78d6"

transacciones <- read_csv("datasets/transacciones.csv",
                          show_col_types = FALSE)
tiendas <- read_csv("datasets/tiendas.csv", show_col_types = FALSE)

tickets_tienda <- transacciones %>%
  left_join(tiendas, by = "tienda_id") %>%
  mutate(tienda = str_remove(nombre_tienda, "Sucursal ")) %>%
  group_by(tienda) %>%
  summarise(tickets  = n(),
            ingresos = sum(total_transaccion),
            .groups  = "drop") %>%
  arrange(desc(tickets))
head(tickets_tienda, 4)

s1 <- ggplot(tickets_tienda,
             aes(x = tickets, y = fct_reorder(tienda, tickets))) +
  geom_col(fill = color_principal, width = 0.7) +
  scale_x_continuous(expand = expansion(mult = c(0, 0.05))) +
  labs(title = "Sur y Plaza A atendieron más tickets en 2023",
       x = "Número de tickets", y = NULL,
       caption = "Fuente: transacciones.csv") +
  theme_minimal() +
  theme(panel.grid.minor = element_blank(),
        panel.grid.major.y = element_blank())
s1
\end{rcode}
\begin{rsalida}
# A tibble: 4 x 3
  tienda  tickets ingresos
  <chr>     <int>    <dbl>
1 Plaza A     109  105216.
2 Sur         109  105305.
3 Express     108  102791.
4 Premium     108  101538.
\end{rsalida}

\textbf{Interpretación.} Sur y Plaza A empatan con 109 tickets, pero la
tienda con más \emph{ingresos} es Outlet (106 tickets): vende un poco menos
de tickets pero de mayor importe. Contar tickets e ingresos responde
preguntas distintas (tráfico contra valor).


\begin{solucion}{m3-2}
Este es el código completo, con su salida y la interpretación.
\end{solucion}

\begin{rcode}
library(tidyverse)
color_principal <- "#2a78d6"
color_resalte   <- "#eb6834"

clientes <- read_csv("datasets/clientes.csv", show_col_types = FALSE)
mediana_edad <- median(clientes$edad)
mediana_edad

s2 <- ggplot(clientes, aes(x = edad)) +
  geom_histogram(binwidth = 5, boundary = 0,
                 fill = color_principal, color = "white") +
  geom_vline(xintercept = mediana_edad, linetype = "dashed",
             color = color_resalte, linewidth = 0.8) +
  annotate("text", x = mediana_edad, y = Inf, vjust = 1.5,
           hjust = -0.1, label = paste("Mediana:", mediana_edad),
           color = color_resalte, size = 3.5) +
  labs(title = "Edad de los clientes",
       subtitle = "Barras de 5 años", x = "Edad", y = "Clientes") +
  theme_minimal() +
  theme(panel.grid.minor = element_blank())
s2
\end{rcode}
\begin{rsalida}
[1] 48.5
\end{rsalida}

\textbf{Interpretación.} La mediana es de 48.5 años y las barras tienen
alturas parecidas entre los 18 y los 75 años: la distribución es más o
menos \emph{pareja} (uniforme), sin una edad dominante. Es lo esperable
porque las edades se simularon al azar en ese rango; en una empresa real
casi siempre verías una concentración en algún grupo de edad.


\begin{solucion}{m3-3}
Los cuatro errores:
\begin{enumerate}
  \item \codigo{fill = "\#2a78d6"} dentro de \codigo{aes()}: el color se
        trata como un dato, sale un color por defecto y una leyenda rara. Va
        fuera de \codigo{aes()}, en la geometría.
  \item El \codigo{+} está al inicio de la línea: R termina la instrucción
        en \codigo{ggplot(...)} y falla con ``Cannot use \codigo{+} with a
        single argument''.
  \item \codigo{geom\_bar()} con \codigo{y}: debe ser \codigo{geom\_col()}
        porque los ingresos ya están calculados.
  \item \codigo{\%>\%} entre capas y \codigo{labels = percent} en un eje de
        pesos: debe ser \codigo{+} y \codigo{labels = dollar}.
\end{enumerate}
(Además conviene agregar \codigo{.groups = "drop"}, aunque aquí no es
obligatorio porque se agrupa por una sola columna.)
\end{solucion}

\begin{rcode}
library(tidyverse)
library(scales)
transacciones <- read_csv("datasets/transacciones.csv",
                          show_col_types = FALSE)

ingresos_pago <- transacciones %>%
  group_by(metodo_pago) %>%
  summarise(ingresos = sum(total_transaccion), .groups = "drop")

s3 <- ggplot(ingresos_pago, aes(x = metodo_pago, y = ingresos)) +
  geom_col(fill = "#2a78d6") +
  scale_y_continuous(labels = dollar)
s3
\end{rcode}

La versión corregida se ejecuta sin errores y muestra cuatro columnas
azules con el eje en pesos. Para dejarla lista para la dirección faltaría
ordenarla, ponerla horizontal y agregar títulos (lista de verificación de
la sección~\ref{sec:m3-antes-despues}).


\begin{solucion}{m3-4}
Este es el código completo, con su salida y la interpretación.
\end{solucion}

\begin{rcode}
library(tidyverse)
color_principal <- "#2a78d6"
color_resalte   <- "#eb6834"
color_gris      <- "#b5b3ad"

transacciones <- read_csv("datasets/transacciones.csv",
                          show_col_types = FALSE)

tickets_mes <- transacciones %>%
  mutate(mes = floor_date(fecha, unit = "month")) %>%
  count(mes, name = "tickets")
promedio <- mean(tickets_mes$tickets)
mes_min  <- tickets_mes %>% slice_min(tickets, n = 1)
promedio
mes_min

s4 <- ggplot(tickets_mes, aes(x = mes, y = tickets)) +
  geom_hline(yintercept = promedio, color = color_gris,
             linetype = "dashed", linewidth = 0.6) +
  geom_line(color = color_principal, linewidth = 0.8) +
  geom_point(color = color_principal, size = 2) +
  annotate("text", x = mes_min$mes, y = mes_min$tickets,
           label = paste("Mínimo:", mes_min$tickets, "tickets"),
           vjust = 1.8, color = color_resalte, size = 3.5) +
  scale_x_date(date_labels = "%b", date_breaks = "1 month") +
  labs(title = "Abril tuvo el menor número de tickets del año",
       subtitle = "Tickets por mes, 2023 (línea punteada = promedio)",
       x = NULL, y = "Tickets") +
  theme_minimal() +
  theme(panel.grid.minor = element_blank())
s4
\end{rcode}
\begin{rsalida}
[1] 83.33333
# A tibble: 1 x 2
  mes        tickets
  <date>       <int>
1 2023-04-01      71
\end{rsalida}

\textbf{Interpretación.} El promedio es de 83.3 tickets por mes y abril
queda muy por debajo con 71. Como abril también fue el peor mes en
ingresos, la caída de ese mes se explica en buena parte por menos tráfico
(menos tickets), no solo por tickets más pequeños.


\begin{solucion}{m3-5}
\codigo{wday(fecha, label = TRUE, week\_start = 1)} devuelve el día como
factor ordenado de lunes a domingo; así el eje queda en el orden correcto
sin hacer nada más.
\end{solucion}

\begin{rcode}
library(tidyverse)
library(scales)
color_resalte <- "#eb6834"
color_gris    <- "#b5b3ad"

transacciones <- read_csv("datasets/transacciones.csv",
                          show_col_types = FALSE)

ingresos_dia <- transacciones %>%
  mutate(dia = wday(fecha, label = TRUE, week_start = 1)) %>%
  group_by(dia) %>%
  summarise(ingresos = sum(total_transaccion), .groups = "drop") %>%
  mutate(destacar = if_else(ingresos == max(ingresos), "Máximo",
                            "Resto"))
ingresos_dia

s_dia <- ggplot(ingresos_dia, aes(x = dia, y = ingresos,
                                  fill = destacar)) +
  geom_col(width = 0.7) +
  geom_text(data = filter(ingresos_dia, destacar == "Máximo"),
            aes(label = dollar(ingresos, accuracy = 1)),
            vjust = -0.5, size = 3.5, color = color_resalte) +
  scale_fill_manual(values = c(Máximo = color_resalte,
                               Resto  = color_gris),
                    guide = "none") +
  scale_y_continuous(labels = dollar,
                     expand = expansion(mult = c(0, 0.1))) +
  labs(title = "El viernes es el día que más vende",
       subtitle = "Ingresos por día de la semana, 2023",
       x = NULL, y = "Ingresos (MXN)") +
  theme_minimal() +
  theme(panel.grid.minor = element_blank(),
        panel.grid.major.x = element_blank())
s_dia
\end{rcode}
\begin{rsalida}
# A tibble: 7 x 3
  dia   ingresos destacar
  <ord>    <dbl> <chr>   
1 lun    128878. Resto   
2 mar    119364. Resto   
3 mié    141329. Resto   
4 jue    145280. Resto   
5 vie    161266. Máximo  
6 sáb    115172. Resto   
7 dom    148703. Resto
\end{rsalida}

\begin{figure}[H]
\centering
\includegraphics[width=0.85\textwidth]{m03-sol-dia}
\caption{Solución del ejercicio m3-5.}
\label{fig:m03-sol-dia}
\end{figure}

\textbf{Interpretación.} El viernes suma \$161,266, contra \$115,172 del
sábado, el día más bajo. Pero recuerda el mapa de calor
(figura~\ref{fig:m03-calor}): el viernes no es el mejor día \emph{todos}
los meses. La diferencia anual existe, pero no es tan consistente como para
basar en ella una estrategia sin más análisis.


\begin{solucion}{m3-6}
Este es el código completo, con su salida y la interpretación.
\end{solucion}

\begin{rcode}
library(tidyverse)
library(scales)
empleados <- read_csv("datasets/empleados.csv", show_col_types = FALSE)

empleados %>%
  group_by(departamento) %>%
  summarise(mediana = median(salario_mensual),
            ric     = IQR(salario_mensual),   # rango intercuartílico
            n       = n(),
            .groups = "drop") %>%
  arrange(desc(mediana))

s6 <- ggplot(empleados,
             aes(x = salario_mensual,
                 y = fct_reorder(departamento, salario_mensual,
                                 .fun = median))) +
  geom_boxplot(fill = "#cde2fb", color = "#2a78d6", width = 0.6) +
  scale_x_continuous(labels = dollar) +
  labs(title = "IT y Operaciones tienen los salarios medianos más altos",
       x = "Salario mensual (MXN)", y = NULL) +
  theme_minimal() +
  theme(panel.grid.minor = element_blank(),
        panel.grid.major.y = element_blank())
s6
\end{rcode}
\begin{rsalida}
# A tibble: 7 x 4
  departamento        mediana    ric     n
  <chr>                 <dbl>  <dbl> <int>
1 IT                   35044. 20725.     9
2 Operaciones          34945. 17114.    10
3 Ventas               32102. 22089.    16
4 Finanzas             31123. 14796.     9
5 Atención al Cliente  30663. 17242.    29
6 Marketing            24193. 11419.    12
7 RRHH                 18291. 13712.    15
\end{rsalida}

\textbf{Interpretación.} IT y Operaciones tienen las medianas más altas
(unos \$35,000) y RRHH la más baja (\$18,291). La mayor dispersión la
tiene Ventas: su rango intercuartílico (\codigo{ric}, el ancho de la caja)
es de \$22,089. Cuidado: algunos departamentos tienen
solo 9 o 10 personas, así que un boxplot con tan pocos datos es
orientativo; siempre acompáñalo del número de observaciones.


\begin{solucion}{m3-7}
Este es el código completo, con su salida y la interpretación.
\end{solucion}

\begin{rcode}
library(tidyverse)
library(scales)
color_principal <- "#2a78d6"
color_resalte   <- "#eb6834"
color_gris      <- "#b5b3ad"

productos <- read_csv("datasets/productos.csv", show_col_types = FALSE)

productos_margen <- productos %>%
  mutate(situacion = if_else(costo > precio_catalogo,
                             "Costo mayor al precio",
                             "Margen positivo"))
count(productos_margen, situacion)

s_costo <- ggplot(productos_margen,
                  aes(x = costo, y = precio_catalogo,
                      color = situacion)) +
  geom_abline(slope = 1, intercept = 0, linetype = "dashed",
              color = color_gris) +
  geom_point(size = 2.5, alpha = 0.8) +
  annotate("text", x = 380, y = 390, label = "precio = costo",
           angle = 20, size = 3, color = "#52514e", vjust = -0.5) +
  scale_color_manual(values = c("Margen positivo" = color_principal,
                                "Costo mayor al precio" = color_resalte)) +
  scale_x_continuous(labels = dollar) +
  scale_y_continuous(labels = dollar) +
  labs(title = "10 de 50 productos cuestan más de lo que se venden",
       subtitle = "Cada punto es un producto del catálogo",
       x = "Costo", y = "Precio de catálogo", color = NULL) +
  theme_minimal() +
  theme(panel.grid.minor = element_blank(),
        legend.position = "top")
s_costo
\end{rcode}
\begin{rsalida}
# A tibble: 2 x 2
  situacion                 n
  <chr>                 <int>
1 Costo mayor al precio    10
2 Margen positivo          40
\end{rsalida}

\begin{figure}[H]
\centering
\includegraphics[width=0.85\textwidth]{m03-sol-costo}
\caption{Solución del ejercicio m3-7: los puntos naranjas, debajo de la
diagonal, tienen margen negativo.}
\label{fig:m03-sol-costo}
\end{figure}

\textbf{Interpretación.} Todo punto por debajo de la diagonal tiene un
costo mayor a su precio: 10 de los 50 productos (20\%). Es un
\emph{hallazgo de calidad de datos} o de negocio que hay que escalar: o el
costo está mal capturado, o se están vendiendo productos con pérdida. Al
área de compras le pedirías revisar primero estos diez productos.


\begin{solucion}{m3-8}
La clave es dibujar dos capas con datos distintos: primero las siete
tiendas en gris (con \codigo{group} para que cada una sea su propia línea y
sin \codigo{color} en \codigo{aes()}, así no entran a la leyenda) y encima
las tres líderes con \codigo{color = tienda}.
\end{solucion}

\begin{rcode}
library(tidyverse)
library(scales)
color_gris   <- "#b5b3ad"
paleta_libro <- c("#2a78d6", "#eb6834", "#1baf7a", "#eda100",
                  "#e87ba4", "#008300", "#4a3aa7", "#e34948")

transacciones <- read_csv("datasets/transacciones.csv",
                          show_col_types = FALSE)
tiendas <- read_csv("datasets/tiendas.csv", show_col_types = FALSE)

acumulado <- transacciones %>%
  left_join(tiendas, by = "tienda_id") %>%
  mutate(tienda = str_remove(nombre_tienda, "Sucursal "),
         mes = floor_date(fecha, unit = "month")) %>%
  group_by(tienda, mes) %>%
  summarise(ingresos = sum(total_transaccion), .groups = "drop") %>%
  group_by(tienda) %>%
  arrange(mes, .by_group = TRUE) %>%
  mutate(acumulado = cumsum(ingresos)) %>%
  ungroup()

# las 3 con más ingresos acumulados en el último mes
top3 <- acumulado %>%
  filter(mes == max(mes)) %>%
  slice_max(acumulado, n = 3) %>%
  pull(tienda)
top3

lideres <- acumulado %>%
  filter(tienda %in% top3) %>%
  mutate(tienda = factor(tienda, levels = top3))  # orden de la leyenda
resto <- acumulado %>% filter(!tienda %in% top3)

s_top3 <- ggplot(mapping = aes(x = mes, y = acumulado)) +
  geom_line(data = resto, aes(group = tienda),
            color = color_gris, linewidth = 0.5) +
  geom_line(data = lideres, aes(color = tienda), linewidth = 0.9) +
  scale_color_manual(values = paleta_libro[1:3]) +
  scale_x_date(date_labels = "%b", date_breaks = "1 month") +
  scale_y_continuous(labels = dollar) +
  labs(title = "Las tres tiendas líderes al cierre de 2023",
       subtitle = "Ingresos acumulados; en gris, las otras 7 tiendas",
       x = NULL, y = "Ingresos acumulados (MXN)", color = NULL) +
  theme_minimal() +
  theme(panel.grid.minor = element_blank(),
        legend.position = "top")
s_top3
\end{rcode}
\begin{rsalida}
[1] "Outlet"  "Sur"     "Plaza A"
\end{rsalida}

\begin{figure}[H]
\centering
\includegraphics[width=0.85\textwidth]{m03-sol-top3}
\caption{Solución del ejercicio m3-8.}
\label{fig:m03-sol-top3}
\end{figure}

\textbf{Interpretación.} Outlet, Sur y Plaza A terminan arriba. Sur y
Plaza A cierran prácticamente empatadas (\$105,305 y \$105,216). El
primer lugar cambió de manos varias veces durante el año y Outlet lo tomó
en noviembre. Como \codigo{top3} se calcula con código, la gráfica se
actualiza sola si cambian los datos.


\begin{solucion}{m3-9}
Dentro de una función, \codigo{\{\{ categoria \}\}} (se lee
``\emph{curly-curly}'') le dice a \paquete{ggplot2} y a \paquete{dplyr} que
use la columna cuyo nombre recibió el argumento.
\end{solucion}

\begin{rcode}
library(tidyverse)
library(scales)
color_principal <- "#2a78d6"

tema_empresa <- function(base_size = 11) {
  theme_minimal(base_size = base_size) +
    theme(plot.title = element_text(face = "bold", color = "#0b0b0b"),
          plot.title.position = "plot",
          axis.text = element_text(color = "#52514e"),
          panel.grid.minor = element_blank(),
          panel.grid.major.y = element_blank(),
          plot.margin = margin(10, 15, 10, 10))
}

grafica_ranking <- function(datos, categoria, valor, titulo) {
  ggplot(datos, aes(x = {{ valor }},
                    y = fct_reorder({{ categoria }}, {{ valor }}))) +
    geom_col(fill = color_principal, width = 0.7) +
    scale_x_continuous(labels = dollar,
                       expand = expansion(mult = c(0, 0.05))) +
    labs(title = titulo, x = NULL, y = NULL) +
    tema_empresa()
}

transacciones <- read_csv("datasets/transacciones.csv",
                          show_col_types = FALSE)
tiendas <- read_csv("datasets/tiendas.csv", show_col_types = FALSE)

por_pago <- transacciones %>%
  group_by(metodo_pago) %>%
  summarise(ingresos = sum(total_transaccion), .groups = "drop")

por_tienda <- transacciones %>%
  left_join(tiendas, by = "tienda_id") %>%
  group_by(nombre_tienda) %>%
  summarise(ingresos = sum(total_transaccion), .groups = "drop")

g_pago <- grafica_ranking(por_pago, metodo_pago, ingresos,
                          "Tarjeta de crédito, el método más usado")
g_tienda <- grafica_ranking(por_tienda, nombre_tienda, ingresos,
                            "Sucursal Outlet lidera en ingresos")

dir.create("resultados", showWarnings = FALSE)
ggsave("resultados/ranking_pago.png", g_pago, width = 7, height = 4,
       dpi = 300, bg = "white")
ggsave("resultados/ranking_tienda.png", g_tienda, width = 7, height = 4,
       dpi = 300, bg = "white")
file.exists(c("resultados/ranking_pago.png",
              "resultados/ranking_tienda.png"))
\end{rcode}
\begin{rsalida}
[1] TRUE TRUE
\end{rsalida}

\textbf{Interpretación.} Con una sola función produces gráficas idénticas
en estilo para cualquier tabla con una columna de categoría y otra de
valor. Es el primer paso hacia un ``kit'' de gráficas del área de BI.


\begin{solucion}{m3-10}
Este es el código completo, con su salida y la interpretación.
\end{solucion}

\begin{rcode}
library(tidyverse)
library(scales)
color_principal <- "#2a78d6"
color_resalte   <- "#eb6834"
color_gris      <- "#b5b3ad"

clientes <- read_csv("datasets/clientes.csv", show_col_types = FALSE)

# (a) KPIs en consola
cat("Clientes totales:", comma(nrow(clientes)), "\n")
cat("% premium:", percent(mean(clientes$es_premium), accuracy = 0.1),
    "\n")
cat("Edad mediana:", median(clientes$edad), "años\n")

# (b) clientes por ciudad, resaltando la mayor
clientes_ciudad <- clientes %>%
  count(ciudad, name = "clientes") %>%
  mutate(destacar = if_else(clientes == max(clientes), "Mayor",
                            "Resto"))

s_ciudades <- ggplot(clientes_ciudad,
                     aes(x = clientes,
                         y = fct_reorder(ciudad, clientes),
                         fill = destacar)) +
  geom_col(width = 0.7) +
  geom_text(aes(label = clientes), hjust = -0.3, size = 3.2,
            color = "#52514e") +
  scale_fill_manual(values = c(Mayor = color_resalte,
                               Resto = color_gris),
                    guide = "none") +
  scale_x_continuous(expand = expansion(mult = c(0, 0.1))) +
  labs(title = "Puebla encabeza por poco: 28 clientes contra 27 de CDMX",
       subtitle = "Clientes registrados por ciudad",
       x = "Clientes", y = NULL,
       caption = "Fuente: clientes.csv") +
  theme_minimal() +
  theme(panel.grid.minor = element_blank(),
        panel.grid.major.y = element_blank())
s_ciudades
\end{rcode}
\begin{rsalida}
Clientes totales: 200 
% premium: 36.5% 
Edad mediana: 48.5 años
\end{rsalida}

\begin{figure}[H]
\centering
\includegraphics[width=0.85\textwidth]{m03-sol-ciudades}
\caption{Solución del ejercicio m3-10, inciso (b).}
\label{fig:m03-sol-ciudades}
\end{figure}

Para el inciso (c), el porcentaje de premium por segmento es el promedio
de la columna lógica \codigo{es\_premium} (\codigo{TRUE} cuenta como 1).
Ordenamos los segmentos por edad, no por valor, porque tienen un orden
natural.

\begin{rcode}
premium_segmento <- clientes %>%
  mutate(segmento = factor(segmento,
                           levels = c("Joven", "Adulto", "Maduro",
                                      "Senior"))) %>%
  group_by(segmento) %>%
  summarise(clientes = n(), pct_premium = mean(es_premium),
            .groups = "drop")
premium_segmento

s_premium <- ggplot(premium_segmento,
                    aes(x = segmento, y = pct_premium)) +
  geom_col(fill = color_principal, width = 0.6) +
  geom_text(aes(label = percent(pct_premium, accuracy = 1)),
            vjust = -0.5, size = 3.5, color = "#52514e") +
  scale_y_continuous(labels = percent, limits = c(0, 0.5),
                     expand = expansion(mult = c(0, 0.05))) +
  labs(title = "El segmento Maduro tiene más clientes premium",
       subtitle = "% de clientes premium por segmento de edad",
       x = NULL, y = NULL, caption = "Fuente: clientes.csv") +
  theme_minimal() +
  theme(panel.grid.minor = element_blank(),
        panel.grid.major.x = element_blank())

dir.create("resultados", showWarnings = FALSE)
ggsave("resultados/clientes_ciudad.png", s_ciudades, width = 7,
       height = 4, dpi = 300, bg = "white")
ggsave("resultados/premium_segmento.png", s_premium, width = 7,
       height = 4, dpi = 300, bg = "white")
\end{rcode}
\begin{rsalida}
# A tibble: 4 x 3
  segmento clientes pct_premium
  <fct>       <int>       <dbl>
1 Joven          21       0.333
2 Adulto         48       0.312
3 Maduro         70       0.414
4 Senior         61       0.361
\end{rsalida}

\textbf{Conclusiones para marketing.}
\begin{enumerate}
  \item Tenemos 200 clientes, 36.5\% premium, con edad mediana de 48.5
        años: la base es madura.
  \item Puebla (28) y CDMX (27) concentran el 27.5\% de los clientes;
        Monterrey es la ciudad con menos (14), una posible oportunidad de
        crecimiento.
  \item El segmento Maduro (41--60 años) es el más grande (70 clientes) y el
        de mayor proporción premium (41\%); los Jóvenes son pocos (21), así
        que su 33\% premium se basa en muy pocos casos.
\end{enumerate}

